% !TeX root = ../main.tex
\begin{center}\textbf{ภาคผนวกประกอบบทความระบบลานิสิต}\end{center}
รายละเอียดข้อกำหนด การประเมิน และสคีมา ณ 18 ตุลาคม 2569
\appendix
\section{เกณฑ์ยอมรับของข้อกำหนดการทำงาน}
เกณฑ์ต่อไปนี้เป็นผลคาดหวังสำหรับทดสอบ ยังไม่ใช่ผลผ่านจริง ควรทดสอบทั้งกรณีอนุญาตและปฏิเสธสิทธิ์โดยเปลี่ยนตัวตนผู้ใช้จริง
\begin{tabularx}{\textwidth}{|p{16mm}|X|}
\hline
รหัส & เกณฑ์ยอมรับ \\ \hline
FR01 & ข้อมูลครบสร้างคำขอ PENDING และขั้นแรก CURRENT; ข้อมูลจำเป็นไม่ครบต้องถูกปฏิเสธโดยไม่เกิดคำขอบางส่วน \\ \hline
FR02 & เจ้าของอ่านใบลาตนเองได้ แต่เข้าถึงของผู้อื่นไม่ได้; แก้หลักฐานได้เฉพาะสถานะที่อนุญาต; เอกสารตรงกับคำขอ \\ \hline
FR03 & ผู้รับผิดชอบขั้นปัจจุบันตัดสินได้; ผู้ที่ไม่เกี่ยวข้องหรือยังไม่ถึงขั้นตัดสินไม่ได้; รายงานจำกัดนิสิตในความดูแล \\ \hline
FR04 & อนุมัติระหว่างทางเปิดขั้นถัดไป; ขั้นสุดท้ายเปลี่ยนเป็น APPROVED; ปฏิเสธเปลี่ยนเป็น REJECTED; รายงานคงขอบเขตหน่วยงาน \\ \hline
FR05 & รับทราบเฉพาะใบลาที่อนุมัติและรายวิชาที่มีความสัมพันธ์; การรับทราบไม่เปลี่ยนผลอนุมัติและไม่เพิ่มยอดใบลา \\ \hline
FR06 & การอ่านรายละเอียดและรายงานต้องตรงกับสิทธิ์เจ้าหน้าที่; เปลี่ยนค่าบทบาทหรือรหัสผู้ใช้ในคำขอแล้วต้องไม่ขยายสิทธิ์ \\ \hline
FR07 & นำเข้าแบบเพิ่มไม่สร้างรายการซ้ำตามคีย์; แบบแทนที่ได้ชุดข้อมูลตามแฟ้ม; รหัสและความสัมพันธ์ถูกต้อง; แถวผิดรูปแบบต้องตรวจพบ \\ \hline
FR08 & ส่งผลซ้ำและตัดสินพร้อมกันได้ผลเดียว; ตัดสินผิดลำดับไม่ได้; แยกแหล่ง MANUAL และกฎอัตโนมัติย้อนหลังได้ \\ \hline
FR09 & โหลดซ้ำไม่เพิ่มยอด; ล้มเหลวไม่เผยแพร่ batch บางส่วน; หนึ่งใบหลายวิชายังนับหนึ่งใบ; เทอมไม่มีใบลามีค่าเป็นศูนย์ \\ \hline
\end{tabularx}
\subsection{การบันทึกหลักฐานและการทดสอบซ้ำ}
หลักฐานของแต่ละกรณีควรประกอบด้วยรุ่นซอร์สโค้ด สภาพแวดล้อม บทบาทผู้ใช้ ข้อมูลตั้งต้น ผลที่คาดหวัง ผลที่พบ และเวลาเก็บข้อมูล หากเกี่ยวข้องกับ warehouse ให้ระบุ batch ที่อ่านและเวลาการโหลดสำเร็จล่าสุดด้วย เพื่อแยกกรณีประมวลผลผิดออกจากกรณีข้อมูลรายงานยังไม่ปรับปรุง
การทดสอบสิทธิ์ควรมีทั้งกรณีที่ได้รับอนุญาตและกรณีที่ต้องถูกปฏิเสธ ผู้ทดสอบต้องเปลี่ยนตัวตนและความสัมพันธ์กับนิสิต ไม่เพียงเปลี่ยนชื่อบทบาทในหน้าจอ เมื่อบริการอ้างอิงไม่พร้อมใช้งาน ควรตรวจว่าระบบตอบความผิดพลาดอย่างเหมาะสมและไม่ขยายขอบเขตข้อมูลโดยอาศัยค่าที่ผู้เรียกส่งเอง
ภาพหน้าจอใช้ยืนยันสิ่งที่ผู้ใช้เห็นได้ แต่ไม่เพียงพอสำหรับยืนยันธุรกรรมหรือการป้องกันการตัดสินซ้ำ จึงควรเก็บผลเรียกบริการและผลตรวจข้อมูลที่จำเป็นควบคู่กัน โดยปิดบังรหัสส่วนบุคคลและข้อมูลลับ การจัดเก็บหลักฐานในรูปแบบเดียวกันช่วยให้ตรวจซ้ำเมื่อแก้กฎหรือเปลี่ยนข้อมูลอ้างอิงได้
\subsection{แผนประเมินประโยชน์และข้อจำกัด}
การประเมินกับผู้ใช้ควรวัดเวลาที่ใช้ยื่นคำขอ ความสำเร็จในการทำงาน และความพึงพอใจ แยกกลุ่มนิสิต อาจารย์ และเจ้าหน้าที่ พร้อมระบุจำนวนผู้เข้าร่วมและเงื่อนไขเปรียบเทียบ หากประเมินบริการควรรายงานเวลาโต้ตอบ เปอร์เซ็นไทล์ และอัตราข้อผิดพลาดภายใต้ภาระงานที่ระบุ บทความฉบับนี้ยังไม่มีข้อมูลดังกล่าวจึงไม่รายงานค่าประสิทธิภาพ
ข้อจำกัดที่เอกสารโครงการระบุคือข้อมูลเจ้าของรายวิชา ภาควิชา และประเภทหลักสูตรยังต้องมี mapping ทางการเพื่อให้รายงานครบถ้วน ค่าที่ไม่ทราบคงไว้เป็นข้อมูลไม่ทราบ ไม่แทนด้วยสัดส่วนจำลอง ประวัติเวลาการเปิดขั้นและตัดสินบางรายการไม่ครบ จึงไม่ควรแทนระยะเวลาที่ขาดด้วยศูนย์ ระบบข้อมูลสนับสนุน เก็บข้อมูลอ้างอิงปัจจุบันโดยไม่มีมิติภาคการศึกษาชัดเจน จึงต้องระวังการผูกข้อมูลนี้กับประวัติย้อนหลัง อีกทั้งยังต้องกำหนดรอบโหลดข้อมูลรายงานอัตโนมัติให้เหมาะกับสภาพแวดล้อมใช้งาน

\clearpage
\section{แผนภาพฐานข้อมูลใบลา (db)}
อ่านสคีมาที่ Deploy เมื่อ 18 ต.ค. 2569: MySQL 9.7.2 จำนวน 19 ตาราง 136 คอลัมน์ รวมตาราง migration เส้นเชื่อมมาจาก foreign key ที่ประกาศจริงเท่านั้น รูปแสดงคีย์เพื่อให้อ่านได้ ส่วนรายการคอลัมน์เต็มอยู่ใน images/db-full.svg และ SQL อยู่ใน db/db/file\_visual.sql ไม่ส่งข้อมูลส่วนบุคคลหรือข้อมูลลับไปยังเครื่องมือออนไลน์
\begin{center}\includegraphics[width=\textwidth,height=0.75\textheight,keepaspectratio]{images/db-keys.png}\end{center}
\clearpage
\section{แผนภาพคลังข้อมูล (wh)}
อ่านสคีมาที่ Deploy เมื่อ 18 ต.ค. 2569: MySQL 9.7.2 จำนวน 19 ตาราง 259 คอลัมน์ ใช้ตระกูล dw\_* โดยข้อเท็จจริงแยกคำขอ รายวิชา และขั้นพิจารณา ไม่ใช้โครงสร้าง wh\_* เดิมแทนฐานที่ติดตั้งจริง คอลัมน์เต็มอยู่ใน images/wh-full.svg และ SQL อยู่ใน db/wh/file\_visual.sql
\begin{center}\includegraphics[width=\textwidth,height=0.75\textheight,keepaspectratio]{images/wh-keys.png}\end{center}
\clearpage
\section{แผนภาพระบบข้อมูลสนับสนุน (leave-com)}
อ่านสคีมาจริงจาก MySQL 8.0.44 ใน Docker บนเครื่องเมื่อ 18 ต.ค. 2569 จำนวน 5 ตาราง 24 คอลัมน์ ประกอบด้วยตารางงาน 4 ตารางและ \_prisma\_migrations อีก 1 ตาราง มี foreign key 2 ความสัมพันธ์ โดย advisor\_student อ้างรหัสจากระบบภายนอก จึงไม่วาด foreign key ไปยังบัญชีในฐานใบลาเอง รายการคอลัมน์เต็มอยู่ใน images/leave-com-full.svg และ SQL อยู่ใน db/leave-com/file\_visual.sql
\begin{center}\includegraphics[width=\textwidth,height=0.65\textheight,keepaspectratio]{images/leave-com-full.png}\end{center}
